\chapter{Gesamtliteraturverzeichnis}
\markboth{Gesamtliteratur}{Gesamtliteratur}

Die sechsundzwanzig Bände nennen zusammen \textbf{siebenhundertachtundzwanzig} kommentierte Quellen. Dieses Verzeichnis wiederholt sie nicht, sondern führt die Kernliteratur auf --- die Standardwerke und Gründungsaufsätze, auf die sich mehrere Bände stützen. Die rechte Spalte nennt die Bände, in deren Quellenverzeichnis das Werk steht.

\begin{merksatz}
\boxtitel{Die zwanzig wichtigsten Bücher}
Wer eine Handbibliothek zur Informatik aufbauen will, beginnt hier.

\begin{center}
{\small
\renewcommand{\arraystretch}{1.34}
\begin{longtable}{@{}P{5.8cm}P{2.8cm}@{}}
\toprule
\textbf{Werk} & \textbf{Gebiet}\\
\midrule
\endhead
Petzold: Code & Einstieg, Band 0\\
Nisan, Schocken: The Elements of Computing Systems & Bit bis Betriebssystem\\
Patterson, Hennessy: Computer Organization and Design & Architektur\\
Cormen u.\,a.: Introduction to Algorithms & Algorithmen\\
Knuth: The Art of Computer Programming & Algorithmen, Referenz\\
Sipser: Introduction to the Theory of Computation & Berechenbarkeit\\
Arora, Barak: Computational Complexity & Komplexität\\
Tanenbaum, Bos: Modern Operating Systems & Betriebssysteme\\
Arpaci-Dusseau: Three Easy Pieces & Betriebssysteme, frei\\
Aho u.\,a.: Compilers (das Drachenbuch) & Übersetzerbau\\
Cover, Thomas: Elements of Information Theory & Information\\
Katz, Lindell: Introduction to Modern Cryptography & Kryptographie\\
Tanenbaum, Wetherall: Computer Networks & Netzwerke\\
Newman: Networks & Graphen\\
Russell, Norvig: Artificial Intelligence & KI\\
Hastie, Tibshirani, Friedman: Elements of Statistical Learning & Lernen\\
Dayan, Abbott: Theoretical Neuroscience & Gehirn\\
Nielsen, Chuang: Quantum Computation & Quanten\-informatik\\
Hofstadter: Gödel, Escher, Bach & Selbstbezug\\
\textbf{Kernighan, Pike: The Practice of Programming} & \textbf{Handwerk}\\
\bottomrule
\end{longtable}}
\end{center}

\begin{center}
\begin{minipage}{0.88\textwidth}
\kursiv{Die letzte Zeile ist gegenüber der vorigen Ausgabe hinzugekommen. Sie steht dort, weil vier der fünf Programmierbände sie unabhängig voneinander an erster Stelle nennen --- und weil sie mit etwa zweihundert Seiten das kürzeste Buch der ganzen Liste ist.}
\end{minipage}
\end{center}
\end{merksatz}
